\documentclass[runningheads]{llncs}
\usepackage[T1]{fontenc}
\usepackage{graphicx}
% Used for displaying a sample figure. If possible, figure files should
% be included in EPS format.
%
% If you use the hyperref package, please uncomment the following two lines
% to display URLs in blue roman font according to Springer's eBook style:
%\usepackage{color}
%\renewcommand\UrlFont{\color{blue}\rmfamily}
%\urlstyle{rm}
%
\begin{document}
%
\title{Incremental Linearization for Integer Arithmetic Revisited}
\author{Marek Dan\v{c}o\orcidID{0009-0008-3031-113X} \and
Karel Chvalovsk\'{y}\orcidID{0000-0002-0541-3889} \and
Mikol\'{a}\v{s} Janota\orcidID{0000-0003-3487-784X}}
% Karel.Chvalovsky@cvut.cz
%
\authorrunning{M. Dan\v{c}o et al.}
% First names are abbreviated in the running head.
% If there are more than two authors, 'et al.' is used.
%
\institute{Czech Technical University in Prague}
%
\maketitle              % typeset the header of the contribution
%
\begin{abstract}
	Incremental Linearization had been previously proposed for solving SMT problems over nonlinear real, and then
	later integer arithmetic constraints. Although conceptually simple, in both cases it had been proven very effective.
	In this paper, we introduce a revised axiom set that achieves faster convergence on exponentiation terms — a class
	of problems where prior axiomatizations struggled. We present a standalone implementation that uses Z3 as a backend
	for linear integer arithmetic, and report on an experimental evaluation over the NIA benchmarks in SMT-LIB.
	Our results show that the approach is competitive with state-of-the-art solvers overall,
	and largely outperforms them on benchmarks involving exponentiation.
	\keywords{SMT \and Nonlinear Arithmetic \and NIA.}

% Incremental Linearization has previously been proposed for solving SMT problems over quantifier-free nonlinear integer arithmetic and has proven effective despite its conceptual simplicity.
%
% In this paper, we introduce a revised axiom set that improves convergence on polynomial constraints, a class of problems on which prior axiomatizations struggled. We present a standalone implementation built on top of Z3 for linear integer arithmetic and evaluate it on the NIA benchmark set from SMT-LIB.
%
% Our results show that the approach is competitive with state-of-the-art solvers overall and substantially outperforms them on benchmarks dominated by polynomial constraints.

\end{abstract}
% \bibliographystyle{splncs04}
% \bibliography{mybibliography}
%
\end{document}
